%%%%%%%%%%%%%%%%%
% CV tailored for Thales - Responsable Offres F/H
%%%%%%%%%%%%%%%%%

\documentclass[8pt,a4paper,ragged2e,withhyper]{altacv}

\geometry{left=1.25cm,right=1.25cm,top=.5cm,bottom=1.cm,columnsep=1.2cm}

\usepackage{paracol}
\usepackage{fontawesome5}
\usepackage{hyperref}

\iftutex 
  \setmainfont{Roboto Slab}
  \setsansfont{Lato}
  \renewcommand{\familydefault}{\sfdefault}
\else
  \usepackage[rm]{roboto}
  \usepackage[defaultsans]{lato}
  \renewcommand{\familydefault}{\sfdefault}
\fi

\definecolor{SlateGrey}{HTML}{2E2E2E}
\definecolor{LightGrey}{HTML}{666666}
\definecolor{DarkPastelRed}{HTML}{8AC0D9}
\definecolor{PastelRed}{HTML}{00B0DA}
\definecolor{GoldenEarth}{HTML}{B4AD0C}

\colorlet{name}{black}
\colorlet{tagline}{PastelRed}
\colorlet{heading}{DarkPastelRed}
\colorlet{headingrule}{GoldenEarth}
\colorlet{subheading}{PastelRed}
\colorlet{accent}{PastelRed}
\colorlet{emphasis}{SlateGrey}
\colorlet{body}{LightGrey}

\definecolor{violet}{HTML}{5A2582}
\definecolor{rouge}{HTML}{F53B3B}
\definecolor{noir}{HTML}{00232C}
\definecolor{bleu}{HTML}{00B0DA}
\definecolor{rose}{HTML}{F831A0}
\definecolor{jaune}{HTML}{FFEC00}
\definecolor{vert}{HTML}{34AD6C}

\renewcommand{\namefont}{\Huge\rmfamily\bfseries}
\renewcommand{\personalinfofont}{\footnotesize}
\renewcommand{\cvsectionfont}{\LARGE\rmfamily\bfseries}
\renewcommand{\cvsubsectionfont}{\large\bfseries}
\renewcommand{\cvItemMarker}{{\small\textbullet}}
\renewcommand{\cvRatingMarker}{\faCircle}

\input{pubs-num.tex}
\addbibresource{sample.bib}

\begin{document}

\name{BOILEAU Guillaume}
\tagline{Technical Bid \& Project Manager | Signal Processing, Complex Systems, Product Strategy \& Cross-Functional Coordination}

\photoL{2.8cm}{moi}

\personalinfo{%
  \email{guillaume.boileaupro@gmail.com}
  \phone{+33 6 59 94 85 84}
  \location{Alpes-Maritimes, FRANCE}
  \linkedin{guillaume-boileau-phd-891b81265}
  \researchgate{Guillaume-Boileau-2}
  \orcid{0000-0002-3576-6968}
}

\makecvheader
\vspace{-0.3cm}

\AtBeginEnvironment{itemize}{\small}
\columnratio{0.64}

\begin{paracol}{2}

\cvsection{Experience}

\cvevent{\textbf{Software Engineer - System Integration, Electrical Equipment Interfaces \& Maintenance Support}}{VEV Platform Services France}{2025 -- 2026}{Valbonne, France}

\textbf{Context:} Development, integration and operational support of software services for EV charging infrastructure within a CPMS platform connected to electrical charging stations and field equipment.

\textbf{Objective:} Support reliable operation of charging infrastructure by ensuring consistency between software services, operational data flows, electrical equipment behaviour and maintenance-related diagnostics.

\begin{itemize}
  \item \textbf{Industrial electrical systems:} Worked on operational services connected to EV charging stations, involving electrical equipment, field infrastructure and maintenance constraints.
  \item \textbf{Maintenance support:} Supported the analysis of issues affecting charging station operation, including equipment status, data exchanges, service behaviour and operational continuity.
  \item \textbf{Electronic/electrical equipment interfaces:} Analysed interactions between CPMS platform services, charging station data, hardware status information and field maintenance needs.
  \item \textbf{System integration:} Contributed to the integration and maintenance of backend and web services interacting with EV charging infrastructure.
  \item \textbf{Operational reliability:} Investigated production issues involving software behaviour, data consistency, electrical equipment feedback and interface continuity.
  \item \textbf{Data flows:} Implemented and monitored FTP-based operational data exchange services used for infrastructure operation and follow-up.
  \item \textbf{Technical coordination:} Interfaced with technical stakeholders to clarify operational problems, understand field constraints and support issue resolution.
  \item \textbf{Agile coordination:} Used Zenhub for task tracking, backlog follow-up, iterative development and technical coordination.
  \item \textbf{Software development:} Developed services using TypeScript, Angular and Node.js in a collaborative delivery environment.
\end{itemize}

\divider

\cvevent{\textbf{Senior R\&D Engineer - Technical Coordination, Validation \& System Studies}}{Laboratoire Lagrange, Observatoire de la Côte d'Azur, CNES}{2023 -- 2025}{Nice, France}

\textbf{Context:} Engineering activities for the LISA ESA/NASA space mission, involving complex software chains, model integration, simulation workflows, validation campaigns and international coordination.

\textbf{Objective:} Coordinate technical activities, structure validation workflows and support mission-level studies under strong performance, reliability, traceability and stakeholder constraints.

\begin{itemize}
  \item \textbf{Complex systems:} Contributed to mission-level analysis workflows for the LISA ESA/NASA space mission within a large international technical collaboration.
  \item \textbf{Technical coordination:} Coordinated contributors, supervised interns and aligned software, modelling and validation activities across multiple stakeholders.
  \item \textbf{Bid-relevant expertise:} Translated complex technical constraints into structured deliverables, validation logic, technical evidence and decision-support material.
  \item \textbf{System studies:} Analysed performance drivers, modelling assumptions, uncertainty sources and system-level limitations affecting mission outputs.
  \item \textbf{Signal processing:} Worked on low-SNR signal analysis, stochastic backgrounds, spectral separation and simulation-based performance assessment.
  \item \textbf{Validation strategy:} Defined comparison campaigns, consistency checks, acceptance logic and review material for technical outputs produced by multiple contributors.
  \item \textbf{Software platform:} Developed CATpyLISA, a Python platform for large-scale model integration, comparison, validation and technical review.
  \textcolor{DarkPastelRed}{\href{https://gitlab.oca.eu/gboileau/lisa_l3_tools}{\bf GitLab:CATpyLISA}}
  \item \textbf{Documentation management:} Used Confluence for technical documentation, coordination notes, validation follow-up and knowledge sharing.
  \item \textbf{Continuous improvement:} Improved workflows, documentation practices and analysis reliability to reduce inconsistencies and strengthen reproducibility.
\end{itemize}

\divider

\cvevent{\textbf{Data Scientist - Signal Analysis, Diagnostics \& Performance Studies}}{Universiteit Antwerpen, Belgium}{2021 -- 2023}{Antwerp, Belgium}

\textbf{Context:} Data analysis and performance studies for precision instruments in a high-requirement technical environment linked to gravitational-wave detector performance.

\textbf{Objective:} Identify, characterize and mitigate sources affecting system sensitivity using robust statistical, computational and validation-oriented methods.

\begin{itemize}
  \item \textbf{Performance studies:} Analysed sources affecting system sensitivity, measurement quality and performance stability.
  \item \textbf{Technical diagnostics:} Identified contributors to performance degradation and supported prioritisation of mitigation actions.
  \item \textbf{Signal and noise analysis:} Developed Python-based pipelines for noise source identification, characterization and mitigation.
  \item \textbf{Decision support:} Produced structured analyses helping technical teams understand constraints, risks and performance limitations.
  \item \textbf{Technical reporting:} Delivered reproducible and traceable technical analyses for international engineering and research stakeholders.
  \item \textbf{System impact:} Contributed to studies identifying environmental and instrumental constraints for next-generation gravitational-wave observatories.
\end{itemize}

\divider

\cvevent{\textbf{Junior R\&D Engineer - Simulation, Signal Processing \& Mission Performance}}{ARTEMIS, Observatoire de la Côte d'Azur}{2018 -- 2021}{Nice, France}

\textbf{Context:} Engineering and analysis activities for gravitational-wave mission preparation, involving signal analysis, simulation, software workflows and noise characterization.

\textbf{Objective:} Develop robust analysis methods for complex signals and support technical investigations in demanding system environments.

\begin{itemize}
  \item \textbf{Signal processing:} Developed methods for the separation of overlapping low-SNR signals in complex datasets.
  \item \textbf{Simulation workflows:} Built simulation approaches for complex signal environments, uncertainty studies and large-scale datasets.
  \item \textbf{Performance analysis:} Supported the identification of limiting factors affecting mission-level performance and system sensitivity.
  \item \textbf{Technical investigations:} Investigated noise sources through cross-sensor correlation analyses in diagnostic contexts.
  \item \textbf{Validation logic:} Compared simulated outputs, model behaviours and analysis results to assess consistency and robustness.
\end{itemize}

\switchcolumn

\cvsection{Profile for Thales}

\textbf{Technical project and bid-oriented engineer with strong experience in signal processing, complex systems, simulation workflows, technical coordination and validation-driven decision support}. Background developed through major international collaborations including \textbf{LISA ESA/NASA}, \textbf{LIGO--Virgo--KAGRA} and \textbf{Einstein Telescope}, complemented by recent industrial experience involving operational services connected to EV charging stations, electrical equipment interfaces and maintenance-related diagnostics.

For a \textbf{Responsable Offres F/H} role within Thales, I bring the ability to understand complex technical subjects, coordinate multidisciplinary contributors, structure technical arguments, clarify risks and constraints, and transform expert inputs into coherent deliverables for decision-making, product strategy and offer construction.

\begin{itemize}
  \item Strong ability to coordinate technical contributors across software, modelling, validation, data analysis and system-level activities.
  \item Experience translating complex scientific and engineering constraints into structured technical documentation and decision-support material.
  \item Background in signal processing, acoustics-related analysis logic, noise modelling, simulation and performance assessment.
  \item Ability to contribute to technical and financial offer construction through structured analysis of scope, risks, assumptions, resources and deliverables.
  \item Experience with roadmap-like thinking: aligning technical developments, validation needs, stakeholder expectations and long-term system performance objectives.
  \item Strong communication skills in English within international collaborations and multidisciplinary technical environments.
  \item Practical experience with documentation, technical reviews, validation evidence, reporting and stakeholder coordination.
  \item Motivation to apply this experience to sonar systems, underwater systems, material products and Thales high-technology industrial environments.
\end{itemize}

\cvsection{Languages}

\cvskill{English}{4}
\divider
\cvskill{Spanish}{2}
\divider

\medskip

\cvsection{Education}

\cvevent{PhD in Physics}{Université Côte d'Azur}{2018 -- 2021}{Nice, France}

\divider

\cvevent{Selective Graduate Program in Fundamental Physics (Magistère)}{Université Paris-Saclay}{2016 -- 2018}{Orsay, France}

\divider

\cvevent{MSc in Astronomy and Astrophysics}{Observatoire de Paris / Université Paris-Saclay}{2017 -- 2018}{Paris, France}

\divider

\cvevent{BSc in Fundamental Physics}{Université Paris-Saclay}{2015 -- 2016}{Orsay, France}

\cvsection{Professional Training}

\cvevent{Supervised Machine Learning (Regression \& Classification)}{DeepLearning.AI - Andrew Ng}{2020}{Coursera}

\end{paracol}

\cvsection{Projects}

\cvevent{CATpyLISA - Technical Coordination, Model Integration \& Validation Platform}{Python Platform for the LISA ESA/NASA Space Mission}{2023 -- 2025}{}

Development of a Python platform dedicated to large-scale model integration, comparison, validation and technical review for LISA mission-level data analysis workflows.

\textbf{Role:} Technical coordination, model integration, software workflow design, validation campaigns, documentation support, consistency checks and interaction with multiple contributors.

\textbf{Relevance for offers and product strategy:} Work involving expert coordination, structured technical deliverables, performance constraints, validation assumptions, risk identification, traceability and decision-support material.

\divider

\cvevent{Co-author and full member of the LISA Consortium}{ESA/NASA Space Mission - Large-scale International Collaboration}{2018 -- now}{}

Contribution to software, analysis and validation workflows for the LISA space mission within an international engineering and scientific collaboration involving ESA, NASA and multiple research institutes.

\textbf{Role:} Software workflows, technical coordination, model integration, simulation campaigns, validation activities, complex signal-processing tasks and collaborative engineering support.

\textbf{Relevance for Thales CCM:} Experience in complex system studies, multidisciplinary collaboration, performance assessment, roadmap-like technical planning and high-standard engineering communication.

\textbf{Program scale:} Major international space mission with an estimated budget of approximately 2 billion EUR over 10 years.

\divider

\cvevent{Co-author and full member of Einstein Telescope and LIGO--Virgo--KAGRA Collaborations}{International Collaborations}{2021 -- now}{}

Contribution to collaborative developments in data analysis, signal characterization, software methods and technical investigations within the LIGO--Virgo--KAGRA and Einstein Telescope collaborations.

\textbf{Role:} Development of analysis tools, contribution to technical activities and support to complex software workflows in high-standard collaborative environments.

\textbf{System impact:} Studies on environmental and instrumental noise sources helping identify fundamental design constraints and performance limitations for next-generation gravitational-wave observatories.

\cvsection{Strengths}

\vspace{-.3cm}

\cvsubsection{Technical Offers, Project Coordination \& Product Strategy}

{\small
\cvtag{Technical Offers}
\cvtag{Bid Support}
\cvtag{Technical Coordination}
\cvtag{Functional Leadership}
\cvtag{Multidisciplinary Teams}
\cvtag{Product Strategy}
\cvtag{Roadmap Logic}
\cvtag{Work Packages}
\cvtag{Scope Analysis}
\cvtag{Risk Analysis}
\cvtag{Assumptions}
\cvtag{Deliverables}
\cvtag{Technical Reviews}
\cvtag{Stakeholder Management}
\cvtag{Decision Support}
\cvtag{Technical Reporting}
}

\divider\smallskip
\vspace{-.3cm}

\cvsubsection{Signal Processing, Systems \& Industrial Equipment}

{\small
\cvtag{Signal Processing}
\cvtag{Acoustics Logic}
\cvtag{Noise Analysis}
\cvtag{Low-SNR Signals}
\cvtag{Simulation Workflows}
\cvtag{Performance Analysis}
\cvtag{System-Level Analysis}
\cvtag{Complex Systems}
\cvtag{Space Systems}
\cvtag{Industrial Equipment}
\cvtag{EV Charging Infrastructure}
\cvtag{Electrical Equipment Interfaces}
\cvtag{Maintenance Support}
\cvtag{Operational Diagnostics}
\cvtag{Model Integration}
\cvtag{Validation Campaigns}
\cvtag{Data Consistency}
\cvtag{Anomaly Analysis}
\cvtag{Root Cause Analysis}
\cvtag{Technical Diagnostics}
\cvtag{Operational Reliability}
}

\divider\smallskip
\vspace{-.3cm}

\cvsubsection{Tools, Software \& Documentation}

{\small
\cvtag{Python}
\cvtag{Linux}
\cvtag{Bash}
\cvtag{Git}
\cvtag{GitLab}
\cvtag{GitHub}
\cvtag{Docker}
\cvtag{CI/CD}
\cvtag{SQL}
\cvtag{TypeScript}
\cvtag{Angular}
\cvtag{Node.js}
\cvtag{REST API}
\cvtag{FTP}
\cvtag{Confluence}
\cvtag{Zenhub}
\cvtag{Agile Tooling}
\cvtag{Technical Documentation}
\cvtag{Technical English}
\cvtag{International Collaboration}
}

\end{document}